\scanpage{87}
\mspage{11}
whereas one will practically never be able to find an $x\in X_K$ \deleted{stable} \inserted{fixed} under $\mathcal G$.
\figure{\begin{tikzpicture}\draw (0,0)--(1.4,0);\end{tikzpicture}}

Before giving a principle of approach for ``symmetric'' calculations, taking into account the presence of a group $\mathcal G$ (a case of great importance in practice), I \begin{tikzpicture}[baseline=(w.base)]\node[inner sep=0pt,outer sep=0pt] (w) {would like};\draw ([xshift=.75em,yshift=-.4em]w.south west) -- ++(-.015,-.05) .. controls +(.03,-.04) and +(-.02,-.02) .. ++(.12,.02) .. controls +(.03,.02) and +(-.04,-.02) .. ++(.10,0) .. controls +(.06,.015) and +(-.05,.015) .. ++(.28,.01) -- ++(.02,-.035);\end{tikzpicture} to make precise the question of the relation between $\Pi_1(X,I)$ ($I$ a finite subset of $X_K$) and $\Pi_1(U_{0,3})$, \deleted{and the} from the viewpoint of the actions of $\Gamma$.
\inserted{To fix ideas I assume that $K$ is \deleted{\unsure{of}} finite over $\mathbb Q$,} \inserted{and} I shall take the case $X$ quasi-projective smooth geometrically connected. Using theorems of Lefschetz type, plus the result of \deleted{Biel} Белый, one \deleted{\unsure{finds|to find}} should be able to find two homomorphisms of schemes over $K$
\diagram{\((44)\)\qquad\begin{tikzcd}[ampersand replacement=\&]
I \arrow[d,phantom,"\cap" description] \& \\
\unsure{\mathcal C|C} \arrow[r,hook,"i"] \arrow[d,"f"'] \& X \\
U_{0,3} \&
\end{tikzcd}}
where $i$ induces a \uline{surjective} homomorphism on the fundamental groups, and where $f$ is
\margin{86}
